\documentclass[11pt]{article}
\usepackage[spanish,es-nodecimaldot]{babel}
\usepackage[utf8]{inputenc}
\usepackage{amsmath,amssymb,geometry,array,tikz}
\usetikzlibrary{babel}
\geometry{letterpaper,margin=1.5cm}
\AtBeginDocument{\setlength{\abovedisplayskip}{3pt}\setlength{\belowdisplayskip}{3pt}\setlength{\abovedisplayshortskip}{2pt}\setlength{\belowdisplayshortskip}{2pt}}
\setlength{\parindent}{0pt}
\renewcommand{\arraystretch}{1.3}
\DeclareMathOperator{\sen}{sen}
\begin{document}
\begin{center}{\Large\textbf{Simulacro 02 --- Matemáticas Aplicadas}}\end{center}
\hrule\vspace{8pt}

%=================== EJERCICIO 1 ===================
{\large\textbf{1) Resolver por Runge-Kutta de 4º}}\\[6pt]
\textbf{Resolver} \quad $y''-3y'=4y+2\sen(t) \qquad y(1)=2 \quad y'(1)=3 \qquad H=0.1$\\[4pt]
\[
\begin{array}{r} f:\\[2pt] g:\end{array}\left\{\begin{array}{l} y'=z\\[2pt] z'=3z+4y+2\sen(t)\end{array}\right.
\qquad\qquad y(1)=2 \qquad z(1)=3
\]
\[
\begin{array}{ll}
K_1=Hf(t,y,z) & C_1=Hg(\,''\,)\\
K_2=Hf\!\left(t+\frac{H}{2},\,y+\frac{K_1}{2},\,z+\frac{C_1}{2}\right) & C_2=Hg(\,''\,)\\
K_3=Hf\!\left(t+\frac{H}{2},\,y+\frac{K_2}{2},\,z+\frac{C_2}{2}\right) & C_3=Hg(\,''\,)\\
K_4=Hf(t+H,\,y+K_3,\,z+C_3) & C_4=Hg(\,''\,)\\[4pt]
\multicolumn{2}{l}{y_{n+1}=y_n+\big(K_1+2(K_2+K_3)+K_4\big)/6 \qquad z_{n+1}=z_n+\big(C_1+2(C_2+C_3)+C_4\big)/6}
\end{array}
\]
{\small
\textbf{Para} $t=1\quad y=2\quad z=3\quad H=0.1$\\[-2pt]
\begin{align*}
K_1&=0.1f(1,2,3)=0.1[3]=0.3\\
C_1&=0.1g(\,''\,)=0.1\left[3(3)+4(2)+2\sen(1)\right]=1.8682\\
K_2&=0.1f\left(1+\tfrac{0.1}{2},\,2+\tfrac{0.3}{2},\,3+\tfrac{1.8682}{2}\right)=0.1f(1.05,\,2.15,\,3.9341)=0.1[3.9341]=0.3934\\
C_2&=0.1g(1.05,\,2.15,\,3.9341)=0.1\left[3(3.9341)+4(2.15)+2\sen(1.05)\right]=2.2137\\
K_3&=0.1f\left(1.05,\,2+\tfrac{0.3934}{2},\,3+\tfrac{2.2137}{2}\right)=0.1f(1.05,\,2.1967,\,4.10685)=0.1[4.10685]=0.4106\\
C_3&=0.1g(\,''\,)=0.1\left[3(4.10685)+4(2.1967)+2\sen(1.05)\right]=2.2842\\
K_4&=0.1f(1+0.1,\,2+0.4106,\,3+2.2842)=0.1f(1.1,\,2.4106,\,5.2842)=0.1[5.2842]=0.5284\\
C_4&=0.1g(\,''\,)=0.1\left[3(5.2842)+4(2.4106)+2\sen(1.1)\right]=2.7277
\end{align*}
\textbf{1\textsuperscript{a} iter}\vspace{-4pt}
\begin{align*}
y(1.1)&=2+\big(0.3+2(0.3934+0.4106)+0.5284\big)/6=2.406\\
z(1.1)&=3+\big(1.8682+2(2.2137+2.2842)+2.7277\big)/6=5.2652
\end{align*}
\textbf{Para} $t=1.1\quad y=2.406\quad z=5.2652\quad H=0.1$\\[-2pt]
\begin{align*}
K_1&=0.1f(1.1,2.406,5.2652)=0.1[5.2652]=0.5265\\
C_1&=0.1g(\,''\,)=0.1\left[3(5.2652)+4(2.406)+2\sen(1.1)\right]=2.7202\\
K_2&=0.1f\left(1.1+\tfrac{0.1}{2},\,2.406+\tfrac{0.5265}{2},\,5.2652+\tfrac{2.7202}{2}\right)=0.1f(1.15,\,2.66925,\,6.6253)=0.1[6.6253]=0.6625\\
C_2&=0.1g(1.15,\,2.66925,\,6.6253)=0.1\left[3(6.6253)+4(2.66925)+2\sen(1.15)\right]=3.2378\\
K_3&=0.1f\left(1.15,\,2.406+\tfrac{0.6625}{2},\,5.2652+\tfrac{3.2378}{2}\right)=0.1f(1.15,\,2.73725,\,6.8841)=0.1[6.8841]=0.6884\\
C_3&=0.1g(\,''\,)=0.1\left[3(6.8841)+4(2.73725)+2\sen(1.15)\right]=3.3426\\
K_4&=0.1f(1.1+0.1,\,2.406+0.6884,\,5.2652+3.3426)=0.1f(1.2,\,3.0944,\,8.6078)=0.1[8.6078]=0.8607\\
C_4&=0.1g(\,''\,)=0.1\left[3(8.6078)+4(3.0944)+2\sen(1.2)\right]=4.0065
\end{align*}
\textbf{2\textsuperscript{a} iter}\vspace{-4pt}
\begin{align*}
y(1.2)&=2.406+\big(0.5265+2(0.6625+0.6884)+0.8607\big)/6=3.0875\\
z(1.2)&=5.2652+\big(2.7202+2(3.2378+3.3426)+4.0065\big)/6=8.5797
\end{align*}
}
\vspace{-6pt}\[ \boxed{\begin{array}{c|ccc} t & 1 & 1.1 & 1.2\\\hline y & 2 & 2.406 & 3.0875\\ z=y' & 3 & 5.2652 & 8.5797\end{array}} \]

{\footnotesize $^{*}$ El enunciado no da el valor de $y'(1)$; se usa $y'(1)=3$ (mismo ejercicio de la hoja anterior).}
\newpage

{\large\textbf{2) Solución de E.D.O de 2º con valores en la frontera usando diferencias finitas}}\\[6pt]
\textbf{Resolver} \quad $y''+4y'+5y=t \qquad y'(1)=2 \quad y(3)=0 \qquad n=5$\\[4pt]
$h=\dfrac{3-1}{5}=\dfrac{2}{5}=0.4$ \hfill
\begin{tikzpicture}[baseline=-3pt,x=1.25cm]
\draw[dashed] (-1,0)--(0,0); \draw (0,0)--(5,0);
\foreach \t/\i in {0/1,1/1.4,2/1.8,3/2.2,4/2.6,5/3}{\draw (\t,0.08)--(\t,-0.08); \node[above,font=\small] at (\t,0.08){\i};}
\foreach \k in {0,...,5}{\node[below,font=\small] at (\k,-0.08){$y_{\k}$};}
\draw (-1,0.08)--(-1,-0.08); \node[below,font=\small] at (-1,-0.08){$y_{-1}$}; \node[above,font=\scriptsize] at (-1,0.08){ficticio};
\end{tikzpicture}
\begin{align*}
&\frac{(y_{i+1}-2y_i+y_{i-1})}{h^2}+4\,\frac{(y_{i+1}-y_{i-1})}{2h}+5y_i=t_i\\[4pt]
&y_{i+1}-2y_i+y_{i-1}+2h(y_{i+1}-y_{i-1})+5h^2y_i=t_ih^2 && (A)\\[4pt]
&(1-2h)\,y_{i-1}-(2-5h^2)\,y_i+(1+2h)\,y_{i+1}=t_ih^2 && (B)
\end{align*}
$y'(1)=\dfrac{y_1-y_{-1}}{2h}=2 \;\Rightarrow\; y_{-1}=y_1-2(2)(0.4)=y_1-1.6 \qquad\qquad y_5=y(3)=0$
\begin{align*}
i=0 \;\; t_0=1 &\quad (1-2(0.4))\,y_{-1}-(2-5(0.4)^2)\,y_0+(1+2(0.4))\,y_1=1(0.4)^2 && (0)\\
&\quad 0.2(y_1-1.6)-1.2\,y_0+1.8\,y_1=0.16 \;\to\; -1.2\,y_0+2\,y_1=0.48\\
i=1 \;\; t_1=1.4 &\quad (1-2(0.4))\,y_0-(2-5(0.4)^2)\,y_1+(1+2(0.4))\,y_2=1.4(0.4)^2 && (1)\\
i=2 \;\; t_2=1.8 &\quad (1-2(0.4))\,y_1-(2-5(0.4)^2)\,y_2+(1+2(0.4))\,y_3=1.8(0.4)^2 && (2)\\
i=3 \;\; t_3=2.2 &\quad (1-2(0.4))\,y_2-(2-5(0.4)^2)\,y_3+(1+2(0.4))\,y_4=2.2(0.4)^2 && (3)\\
i=4 \;\; t_4=2.6 &\quad (1-2(0.4))\,y_3-(2-5(0.4)^2)\,y_4+(1+2(0.4))\,y_5=2.6(0.4)^2 && (4)\\
&\quad 0.2\,y_3-1.2\,y_4+1.8(0)=0.416
\end{align*}
\[
\underbrace{\begin{bmatrix}
-1.2 & 2 & 0 & 0 & 0\\
0.2 & -1.2 & 1.8 & 0 & 0\\
0 & 0.2 & -1.2 & 1.8 & 0\\
0 & 0 & 0.2 & -1.2 & 1.8\\
0 & 0 & 0 & 0.2 & -1.2
\end{bmatrix}}_{A}
\underbrace{\begin{bmatrix} y_0\\ y_1\\ y_2\\ y_3\\ y_4\end{bmatrix}}_{X}
=
\underbrace{\begin{bmatrix} 0.48\\ 0.224\\ 0.288\\ 0.352\\ 0.416\end{bmatrix}}_{B}
\qquad X=A^{-1}B
\]
\[
\boxed{\begin{array}{c|cccccc}
t & 1 & 1.4 & 1.8 & 2.2 & 2.6 & 3\\\hline
y & -21.776 & -12.8256 & -6.0064 & -2.4192 & -0.7498 & 0
\end{array}}
\]

\newpage

{\large\textbf{3) Teoría de la variable compleja --- Residuos}}\\[6pt]
\[
\textbf{Evalúe}\;\int\limits_{C}\frac{2z^2+5}{(z^2+2z+1)(z+1)}\,dz=2\pi i\sum_{i=1}^{n}R_i=2\pi i(2)=4\pi i
\]
\begin{minipage}{0.55\textwidth}
$C: |z+1|=2$\\[6pt]
$z_0=-1 \qquad r=2$\\[6pt]
$z^2+2z+1=(z+1)^2 \;\to\; (z+1)^2(z+1)=(z+1)^3$\\[6pt]
$(z+1)^3=0 \;\to\; z=-1 \quad$ (polo de orden 3)
\end{minipage}\hfill
\begin{minipage}{0.4\textwidth}
\begin{tikzpicture}[scale=0.75]
\draw[->] (-3.6,0)--(1.8,0) node[right]{$x$}; \draw[->] (0,-2.5)--(0,2.6) node[above]{$i$};
\draw[thick] (-1,0) circle (2);
\draw[->] (-1,0)--({-1+2*cos(45)},{2*sin(45)}) node[midway,above left]{\small $R=2$};
\node at (-1,0){$\times$}; \node[below] at (-1,-0.1){\small $-1$};
\end{tikzpicture}
\end{minipage}\\[6pt]
$F(z)=\dfrac{2z^2+5}{(z+1)^3}$
\[
R=\frac{1}{(3-1)!}\lim_{z\to-1}\frac{d^2}{dz^2}\{(z+1)^3F(z)\}
=\frac12\,\frac{d^2}{dz^2}\left\{\begin{matrix}2z^2+5\\ 4z\\ 4\end{matrix}\right\}\Bigg|_{z=-1}=\frac12(4)=2
\]
\[
\boxed{\int\limits_{C}\frac{2z^2+5}{(z^2+2z+1)(z+1)}\,dz=2\pi i(2)=4\pi i\approx 12.5663\,i}
\]
{\footnotesize $^{*}$ Si el denominador fuera $(z^2+2z+5)(z+1)$, los polos $-1\pm2i$ caerían \emph{sobre} $C$ ($|\pm2i|=2$) y la integral no estaría definida.}

\end{document}
